\documentclass[11pt]{article}

% ---------------------------------------------------------------------------
%  Foundational_Amplitudes -- definitions, methods & results
%  A running, mathematically precise record of the amplitude foundation-model
%  program: the regression target and its preprocessing, the equivariant model
%  and its physical encodings, the multi-fidelity HPO machinery, and the
%  empirical results -- optimization laws, feature ablations, scaling, and the
%  training-throughput campaign -- each stated so a reader fluent in physics and
%  ML can reproduce and audit it.
% ---------------------------------------------------------------------------

\usepackage[margin=1in]{geometry}
\usepackage{amsmath, amssymb, amsthm, mathtools}
\usepackage{booktabs}
\usepackage{multirow}
\usepackage{bm}
\usepackage{graphicx}
\usepackage{float}          % the [H] specifier
\graphicspath{{figs/}{../analysis/hpo_optima/}{../compare_models/figs/}{../plots/}{../analysis/divergences/figs/}}
\usepackage{xcolor}
\usepackage[colorlinks=true, linkcolor=blue!55!black, citecolor=blue!55!black]{hyperref}

% --- notation -------------------------------------------------------------
\newcommand{\R}{\mathbb{R}}
\newcommand{\E}{\mathbb{E}}
\newcommand{\Var}{\operatorname{Var}}
\newcommand{\norm}[1]{\left\lVert #1 \right\rVert}
\newcommand{\abs}[1]{\left\lvert #1 \right\rvert}
\newcommand{\Mamp}{\mathcal{M}}
\newcommand{\amp}{\abs{\Mamp}^2}
\newcommand{\Lloss}{\mathcal{L}}
\newcommand{\vlnr}{\mathcal{L}^{\mathrm{val}}_{0}}   % val loss, no reg
\newcommand{\lrpeak}{\eta_{\mathrm{peak}}}
\newcommand{\tstar}{t_\ast}
\newcommand{\nh}{n_h}
\newcommand{\code}[1]{\texttt{#1}}

\newtheorem{definition}{Definition}
\newtheorem{remark}{Remark}
\newtheorem{principle}{Principle}

\title{\textbf{A Foundation Model for Scattering Amplitudes}\\[2pt]
\large Definitions, Methods, and Results}
\author{Joaqu\'in Iturriza}
\date{\today}

\begin{document}
\maketitle

\begin{abstract}
We train a single Lorentz-equivariant transformer to regress tree- and loop-level
scattering amplitudes across many processes ($ee\!\to\!WW$, $ee\!\to\!t\bar t$,
$ee\!\to\!u\bar u(g)$, $ee\!\to\!\gamma\gamma(\gamma)$, \dots) jointly, as a step
toward a reusable ``foundation'' model that transfers to new processes and
perturbative orders. This note fixes precise definitions for the regression target
and its preprocessing, the equivariant architecture and the physical particle
encoding, the joint loss, and the multi-fidelity (\textsc{DyHPO}) hyperparameter
search, and then reports the empirical results with their governing quantities. The
settled results are: (i) amplitude preprocessing must be resolved \emph{per process},
which lowers the median relative error by $\sim\!6\times$ and de-sensitizes the whole
sweep; (ii) under maximal-update parameterization ($\mu$P) the optimal learning rate
\emph{transfers across model width} but obeys a strong \emph{inverted-$U$ law in
training length}, peaking at an absolute optimizer time scale $\tstar\!\approx\!3\times10^3$
steps independent of dataset size --- which rewrites how we set and search learning
rates; (iii) a family of physics-motivated input features (coupling order, internal
propagator mass, off-shellness, \dots) helps exactly when the added quantity is
\emph{not recoverable from the external momenta}, and is neutral otherwise; (iv) among
the three maintained $\mu$P architectures LLoCa is both the cheapest ($\sim\!20\times$ fewer
FLOPs/step than L-GATr) and the most accurate at matched compute \emph{and} matched wall-time,
and fine-tuning a foundation model onto held-out one-loop targets transfers with the same
scaling exponent but a strictly lower floor and no plateau, with full per-layer-decay retraining
the best method; and (v) a per-step optimization campaign took the training step from
$0.356$ to $0.128$\,s (data-bound fix) down onto a compute floor that vectorization and
a fused optimizer lowered by a further $\sim\!3\times$, every numerical fast path proven
equivalent to the original. We close with the methodological principles these results
impose on how sweeps and HPO are run.
\end{abstract}

\tableofcontents

% ==========================================================================
\section{The system, target, and notation}\label{sec:system}
% ==========================================================================

\paragraph{Scattering amplitudes.} For a fixed partonic process the squared matrix
element $\amp$ is a function of the external four-momenta $\{p_i\}_{i=1}^{n}$ (with
$p_i\in\R^{4}$, metric $(+,-,-,-)$) that determines the differential cross section. It
is the object a Monte-Carlo event generator evaluates once per phase-space point; our
aim is to \emph{learn} it, so that a trained network replaces the per-point matrix-element
call. We treat the problem as regression of a single scalar, $\amp$, from the event's
particle set.

\paragraph{Data.} Each dataset is one process, stored as rows of flat four-momenta
($4n$ numbers), the $n$ PDG identifiers, and the scalar $\amp$ in the last column.
Events are generated over a range of centre-of-mass energy $\sqrt s$. The maintained
joint set spans eight base processes ($ee\!\to\!WWZ,\,WW,\,t\bar t,\,u\bar u,\,u\bar ug,\,
u\bar ugg,\,\gamma\gamma,\,\gamma\gamma\gamma,\,u\bar u$) and, in the ``big-run'' recipes,
of order $10^{2}$ datasets once the coupling-, mass-, and order-scans are included. A
\emph{process} (or \emph{dataset}) index $p\in\{1,\dots,P\}$ labels which physical
process an event belongs to; the model is trained on all of them at once.

\paragraph{Notation.} We write $n$ for the number of particles in an event, $\nh$ for the
number of attention heads (the $\mu$P \emph{width} axis, with transformer hidden dimension
$d=16\,\nh$), $D\equiv n_{\mathrm{train}}$ for the number of training events, and
$t\equiv t_{\mathrm{steps}}$ for the number of optimizer steps (the training horizon /
\emph{fidelity}). $\eta$ denotes the learning rate \code{training.lr}. A hat marks a
model prediction. All logarithms in scaling exponents are base-10 unless noted.

% ==========================================================================
\section{The regression target and its preprocessing}\label{sec:target}
% ==========================================================================

Raw $\amp$ spans many orders of magnitude (down to $\sim\!10^{-13}$) and can be
negative at one loop (interference); it must be compressed and standardized before it
is a sane regression target.

\begin{definition}[Amplitude transform]\label{def:trafo}
The target is the preprocessed amplitude $y = (s\circ f)(\amp)$, a \emph{log} (or
\emph{signed-log}) compression followed by an affine \emph{standardization}:
\begin{equation}
  f(a) =
  \begin{cases}
    \ln a, & \text{$a>0$ throughout (positive-definite $\amp$),}\\[2pt]
    \operatorname{sgn}(a)\,\ln\!\bigl(1+\abs a\bigr), & \text{$a$ genuinely signed (loop interference),}
  \end{cases}
  \qquad
  s(x) = \frac{x-\mu}{\sigma},
  \label{eq:trafo}
\end{equation}
with $(\mu,\sigma)$ estimated on the training split only. The choice log vs.\ signed-log
is made by \code{resolve\_amp\_trafos} from the sign content of the data.
\end{definition}

\begin{principle}[Per-process preprocessing]\label{prin:perproc}
The transform and its $(\mu,\sigma)$ are resolved \textbf{per process}, not once
globally. This is the validated default (\code{preprocess\_per\_dataset: true}).
\end{principle}

\noindent The reason is quantitative and was a genuine bug fix:
signed-log only compresses $\abs a\!\gg\!1$, but every $\amp\!\ll\!1$, so a single
\emph{global} signed-log --- forced on all processes by one loop-level dataset with
negative entries --- leaves the targets in \emph{linear} scale (target $\sigma$ ranging
$8$--$30$, up to $1.2\times10^{4}$ for one process), whereupon a shared standardization
cannot reach unit variance and most processes simply fail to train. Resolving the
transform per process ($416$ log / $32$ signed-log across the $448$-dataset big run)
collapses the global preprocessed $\sigma$ from $282$ to $1.06$.

% ==========================================================================
\section{The model and the physical encoding}\label{sec:model}
% ==========================================================================

\paragraph{Architecture.} The maintained models are $\mu$P (maximal-update
parameterized) Lorentz-equivariant transformers, and the default is the \textbf{LLoCa}
(Lorentz Local Canonicalization) transformer (\code{LLOCAMuPTransformer}).

\begin{definition}[Lorentz Local Canonicalization]\label{def:lloca}
LLoCa (arXiv:2505.20280) makes an \emph{arbitrary} backbone exactly Lorentz-equivariant by
canonicalizing into equivariantly-predicted \emph{local frames} rather than by building
equivariance into every layer. A small frames network (\code{LearnedPDFrames}) reads the
event and outputs, for each particle $i$, a Lorentz transformation $\Lambda_i\in SO^+(1,3)$
that varies equivariantly with the input (if the event is boosted/rotated by $\Lambda$, then
$\Lambda_i\!\to\!\Lambda_i\Lambda^{-1}$). Each particle's four-vector and tensor features are
mapped into its own frame, $x_i\mapsto \Lambda_i x_i$, which renders them \emph{Lorentz
invariant}; an ordinary (non-equivariant) backbone --- here a transformer --- then processes
these invariant local features, and any tensorial output is mapped back through $\Lambda_i^{-1}$,
so the network as a whole is \emph{exactly} equivariant. Because the backbone is a plain
transformer, LLoCa is $\sim\!4\times$ faster at $\sim\!10\times$ fewer FLOPs than specialized
equivariant layers at matched accuracy (Table~\ref{tab:arch}), which is why it is the default
here rather than the fully-tensorial L-GATr.
\end{definition}

\noindent The backbone (\code{MuPTransformer}) attends over the event's particle set with a
block-diagonal mask built from the event boundaries \code{ptr}, so particles attend only within
their own event (variable-length events are packed into one flat batch). Two fully-tensorial
$\mu$P L-GATr variants (\code{lgatr\_mup}, \code{lgatr\_slim}) are maintained alternatives.
Defaults: $8$ blocks, $\nh=8$ heads, \code{attn\_reps}$=$``\code{8x0n+2x1n}'' (the frames carry
$8$ scalar and $2$ vector representations through attention).

\begin{definition}[Maximal-update parameterization ($\mu$P)]\label{def:mup}
$\mu$P is a width-dependent rescaling of initialization and per-layer learning rates
under which the optimal hyperparameters --- above all $\eta$ --- become
\emph{width-invariant}, so a configuration tuned at a small width can be reused at any
larger width without re-tuning. The width axis here is $\nh$ (base $\nh=2$, delta $\nh=8$
for the base-shape file); \code{attn\_reps} is held fixed between base and delta.
\end{definition}

\begin{definition}[Physical particle encoding]\label{def:enc}
With \code{use\_PIDs: false} (the default) each PDG id maps through a fixed global table
to an $8$-dimensional \emph{physical property vector}
\[
  \bigl[\,Q,\ \text{spin},\ \log_{10}(m/\mathrm{GeV}),\ T_3,\ B,\ L,\ \text{color charge},\ C_2(R)\,\bigr]
\]
(electric charge, spin, log-mass, weak isospin, baryon and lepton number, and the SU(3)
color charge and quadratic Casimir). The two \emph{categorical} columns --- spin and the SU(3)
color representation --- are \textbf{one-hot} encoded rather than passed as scalars, because
they are labels, not magnitudes (spin-$1$ is not ``twice'' spin-$\tfrac12$), so a scalar column
would force every value onto one line through the origin; the continuous columns are
standardized, and an explicit ``exactly massless'' flag neutralizes the log-mass sentinel. A
learned encoder --- a $2$-layer MLP ($\operatorname{Linear}\!\to\!\text{GELU}\!\to\!\operatorname{Linear}$,
so it can form nonlinear products of quantum numbers) by default --- projects the resulting
$n_{\mathrm{feat}}$-vector to the $16$-D particle embedding. Because the input dimension depends
only on the number of \emph{properties}, never on the vocabulary, adding a new particle or a
$9$th quantum number never forces retraining the transformer. The legacy \code{use\_PIDs: true}
uses a one-hot token index instead. (One-hot color and the MLP embedder are on by default; their
ablation is \S\ref{sec:arch}.)
\end{definition}

\begin{definition}[Diagram graph conditioning]\label{def:diag}
Beyond per-particle properties, a graph encoder embeds each process's set of tree-level
Feynman diagrams into a vector $E(\text{process})$ (a Lorentz scalar), broadcast to every
particle and concatenated into the scalar features exactly like the coupling labels below.
This hands the model the process's diagrammatic/topological structure (which propagators
connect which legs) that is not otherwise recoverable from a permutation-blind particle set.
It is on by default (\code{use\_diagrams: true}); processes with a missing diagram sidecar
degrade gracefully to a zero embedding.
\end{definition}

\begin{definition}[Coupling / auxiliary scalars]\label{def:coup}
Per-dataset physics that is \emph{not} a function of the external momenta is broadcast to
every particle as extra scalar features: the coupling order
\code{amp\_orders}$=[n_{\mathrm{loops}},\,\alpha_s\text{-power}]$
(LO $=[0,0]$, NLO-full $=[1,1]$, virtual-only $=[1,0]$), the per-event log-coupling
$\log\alpha$, and, for a scanned propagator/resonance, the internal log-mass
$\log(m/m_{\mathrm{table}})$. Mixing orders or couplings therefore needs no architecture
change. These are the ``feature levers'' ablated in \S\ref{sec:levers}.
\end{definition}

% ==========================================================================
\section{The joint loss and the comparison metric}\label{sec:loss}
% ==========================================================================

\begin{definition}[Per-process joint loss]\label{def:loss}
For a batch mixing $P$ processes, let $m_p=\frac{1}{\abs{B_p}}\sum_{e\in B_p}(\hat y_e-y_e)^2$
be process $p$'s mean squared error on the preprocessed target ($B_p$ its share of the batch).
The training loss is the \textbf{geometric mean} over processes --- a mean in log-space --- plus
regularization,
\begin{equation}
  \Lloss \;=\; \exp\!\Bigl(\tfrac1P\textstyle\sum_{p=1}^{P}\ln m_p\Bigr)
   \;+\; \lambda\,R(\theta),
  \label{eq:loss}
\end{equation}
with $R(\theta)$ an $L_2$/$L_1$ weight penalty and $\lambda=$\code{regularization\_lambda}.
\end{definition}

\begin{principle}[Aggregate the joint loss in log-space]\label{prin:geomean}
The geometric mean (\code{loss\_aggregation: geometric\_mean}), not the arithmetic mean, is the
default and is used in essentially every run. Processes of different final-state dimensionality
plateau at very different MSE magnitudes (Table~\ref{tab:scaling}: e.g.\ $\sim\!10^{-6}$ for a
$2\!\to\!2$ process vs.\ $\sim\!10^{-2}$ for a $2\!\to\!4$ one). An arithmetic mean is then
dominated by the few high-MSE processes and their gradient, freezing the already-well-learned
ones; the geometric mean weights each process's \emph{relative} improvement equally --- a $10\%$
drop on a $10^{-6}$ process counts the same as a $10\%$ drop on a $10^{-2}$ one --- so every
process keeps scaling regardless of its absolute loss.
\end{principle}

\paragraph{The training sampler: balanced vs.\ equal.} A batch's per-process composition is set
by the sampler. The default is the \textbf{naive equal (uniform) sampler}: draw uniformly over
all events, i.e.\ each process in proportion to its dataset size. A \code{ProcessBalancedSampler}
was built as an alternative that dynamically re-weights processes toward those still improving:
each validation it fits a local scaling exponent $\alpha_p=-\mathrm{d}\log m_p/\mathrm{d}\log
C_p$ (loss vs.\ compute-spent-on-$p$) with a robust slope and its standard error, and --- only on
a significance-gated call ($\alpha_p+k\,\mathrm{SE}<\alpha_{\min}$, failing safe to ``keep
feeding'' when noisy) --- throttles a confirmed-plateaued process and boosts those scaling below
their expected solo rate. It is a genuinely interesting controller, but a controlled A/B (8
processes, best-of-sweep with the controller's own knobs exposed to DyHPO) found the equal
sampler \emph{wins} best-vs-best on $\vlnr$: $8.1\times10^{-8}$ (equal) vs.\ $2.6\times10^{-7}$
(balanced), a $\sim\!3\times$ gap. The benefit dilutes with many datasets and perturbing the
sampling dynamics costs more than it gains, so the balanced sampler was \textbf{discarded as the
default} (\code{use\_balanced\_sampler: false}) and kept opt-in for mixtures of a \emph{few, very
unbalanced} datasets, where the mechanism may still pay.

\begin{definition}[Comparison metric $\vlnr$]\label{def:metric}
Every model comparison in this note uses the \textbf{best non-regularized validation loss}
on the log-amplitude target,
\begin{equation}
  \vlnr \;=\; \min_{t\le T}\ \Lloss_{\text{val}}(t)\big|_{\lambda=0},
  \label{eq:vlnr}
\end{equation}
i.e.\ Eq.~\eqref{eq:loss} \emph{without} the $\lambda R$ term, evaluated on the validation
split and minimized over the run. It is the object DyHPO optimizes and the checkpoint
selector tracks. Using $\vlnr$ (not the regularized tracked loss) is essential: $\lambda$
is itself a tuned hyperparameter, so comparing regularized losses would confound the metric
with the search. Two sides are only comparable when they share the preprocessing of
\S\ref{sec:target}; a feature that changes standardization changes the scale of $\vlnr$
and must be re-based before comparison.
\end{definition}

% ==========================================================================
\section{Multi-fidelity hyperparameter search (DyHPO)}\label{sec:dyhpo}
% ==========================================================================

\begin{definition}[Fidelity and DyHPO]\label{def:dyhpo}
A \emph{fidelity} is a cheap-to-expensive budget knob --- most often the training-step count
$t$, but in principle anything (dataset size, width) --- along which a configuration's
performance can be extrapolated. DyHPO is a multi-fidelity Bayesian optimizer built on a
\emph{deep-kernel} Gaussian process: an MLP feature extractor consumes the HP vector
\emph{concatenated with the (normalized) budget}, and its output feeds an exact GP; a
permutation-invariant set encoder (DeepSets) summarizes a configuration's already-observed
$(\text{budget},\text{loss})$ pairs. Because the budget is a continuous \emph{input} to the
kernel (not a bucketing label), the surrogate \emph{extrapolates} a partially-trained
configuration's loss to higher, unobserved fidelities, and promotes the configuration whose
extrapolated Expected Improvement wins. All jobs of a sweep share one surrogate through a lock
file on the shared filesystem, so a trial informs its siblings only if it \code{observe}s
before they \code{suggest}.
\end{definition}

\begin{remark}[Multi-axis fidelity, and the single-fidelity limit]\label{rem:dyhpo}
The implementation here generalizes DyHPO from a single scalar budget to a
\textbf{multi-dimensional fidelity space}: an independent budget axis per dataset plus $t$, with
the candidate fidelities the Cartesian product of the per-axis levels, each axis normalized and
fed as its own continuous kernel input (this replaces the upstream fixed-length CNN
learning-curve encoder with the DeepSets set encoder, since there is no single ordered curve
across a grid). It reduces cleanly to ordinary single-fidelity Bayesian optimization when only
one budget value is passed: the budget input is then constant, the DeepSets context empty, and
the surrogate collapses to a GP over the HP space with Expected-Improvement acquisition. This
single-fidelity, plain-Bayesian mode is how nearly every sweep in the project is actually run ---
which makes the submission ordering of \S\ref{sec:principles} (interleave, or fall back to
sequential batches) the thing that determines whether the GP has anything to fit.
\end{remark}

A \emph{sweep} draws HP candidates from a declared \emph{search space} (log-uniform or
uniform ranges per hyperparameter) and runs a fixed budget of trials across the fidelity
schedule. The results of \S\ref{sec:lrlaw} are what those search spaces should be; the
methodological cost of getting them wrong is the subject of \S\ref{sec:principles}.

% ==========================================================================
\section{Optimization laws: how the learning rate should be set}\label{sec:lrlaw}
% ==========================================================================

This section is the one that changes how we run sweeps. It is harvested from
\textbf{771} DyHPO sweeps on disk, of which \textbf{422} are \emph{converged} in the
operational sense that their effective learning-rate knob ($\eta$ for pretrain/solo/scaling
sweeps, \code{fine\_tune.lr\_scale} for finetune) does not sit at the edge of its explored
range and the sweep has $\ge 6$ observations ($338$ $\eta$-sweeps, $84$
\code{lr\_scale}-sweeps). Each sweep contributes its own best $\vlnr$ and the HP
configuration attaining it.

% --------------------------------------------------------------------------
\subsection{$\mu$P holds across width; it does not excuse a fixed $\eta$ across horizon}
% --------------------------------------------------------------------------

The $\mu$P promise (Def.~\ref{def:mup}) is that $\eta$ transfers across width. It does:
binning the best $\eta$ by width and taking the geometric mean gives, across a $16\times$
range $\nh\in\{2,\dots,32\}$,
\[
  \eta^\ast(\nh):\quad 6.0,\ 3.6,\ 3.4,\ 3.0,\ 2.9\ \ (\times10^{-3}),
\]
i.e.\ \textbf{flat} --- re-sweeping $\eta$ per width is pure waste (Fig.~\ref{fig:hpo}, top-left).
But $\eta^\ast$ is \emph{not} flat in the training horizon $t$: it rises, then falls.

\begin{principle}[Inverted-$U$ learning-rate law]\label{prin:invU}
At fixed data, the optimal learning rate rises then falls with the training horizon,
peaking at an \textbf{absolute optimizer time scale} $\tstar\approx 3\times10^{3}$ steps:
\begin{equation}
  \eta^\ast(t) \;\approx\; \lrpeak\,\min\!\Bigl[(t/\tstar)^{+0.5},\ (t/\tstar)^{-0.55}\Bigr],
  \qquad \lrpeak \approx 1\times10^{-2},\ \ \tstar\approx 3\times10^{3}.
  \label{eq:invU}
\end{equation}
\end{principle}

\noindent Equation~\eqref{eq:invU} is an empirical approximation --- a convenient form for
\emph{centering} the learning rate at a given horizon, not a claim that either branch is an exact
power law. Each side is roughly log-linear and consistent across the joint ($8$-process,
$\mathrm{bs}=256$) and solo ($\mathrm{bs}=16384$) regimes: ascending slope
$\mathrm{d}\log\eta/\mathrm{d}\log t \approx +0.43$ to $+0.51$ ($t\lesssim\tstar$), descending
slope $-0.55$ to $-0.58$ ($t\gtrsim\tstar$). The geometric-mean $\eta^\ast$ by horizon is
$7.8\times10^{-4}$ at $t=4$, $1.1\times10^{-3}$ at $t=100$, $6$--$8\times10^{-3}$ at the
$1$--$3$k peak, then $2.5\times10^{-3}$ ($t=10$k), $1.6\times10^{-3}$ ($t=31$k),
$7\times10^{-4}$ ($t=100$k) (Fig.~\ref{fig:iters}).

\begin{remark}[Why an absolute step scale, not the convergence point]
The peak sits at $\operatorname{argmax}_t = 3162$ steps for \emph{every} dataset size from
$700$ to $70{,}000$ (a $100\times$ range), whereas the loss-floor knee shifts \emph{later}
with data. So $\tstar$ is not ``where the model converges'' --- it is an optimizer/schedule
time scale: Adam's second-moment EMA has horizon $1/(1-\beta_2)\approx10^{3}$ steps and warmup
is a fraction of $t$ (the optional weight-EMA is \emph{not} implicated --- it was inactive by
default in every sweep these optima are harvested from). Below
$\tstar$ the run never leaves warmup/steady-state transient, so a large $\eta$ cannot be
exploited and the optimum is \emph{lower} and rising (the counter-intuitive branch we rarely
sweep); above $\tstar$ extra steps polish an already-fit model and $\eta^\ast$ falls as
$t^{-0.55}$ (the intuitive branch).
\end{remark}

\begin{remark}[The optimum is a 2D surface $\eta^\ast(t,D)$: evaluate it at your own $(t,D)$]
Disentangled on the clean 2D grid (Fig.~\ref{fig:2d}), the optimum rises \emph{weakly} with
data, $\eta^\ast\propto D^{+0.15\ldots0.3}$, and \emph{saturates} (measured $\sim\!2.2\times$ from
$D=700$ to $70$k at $t=3162$, most of it by $D\!\sim\!2$k). The optimum is thus a 2D surface
$\eta^\ast(t,D)$ --- the inverted-$U$ in $t$ (Eq.~\eqref{eq:invU}) times this weak, saturating
factor in $D$ --- and both axes are set by the experimenter when the run is designed.
\textbf{Set the learning rate by evaluating the surface at your own $t$ and $D$}, sweeping
$\pm\tfrac12$ decade around that value --- not the $t$-only $\eta_c(t)$, which discards the
measured $D$-axis. A large-$D$ run sits at the high-$D$ asymptote, so its centre lies
\emph{above} the $D$-pooled $\eta_c(t)$; read it off the high-$D$ row of the grid at your $t$.
Two cautions the pooled curve hides: (i) the high-$D$ rows only reach $t\approx1.8\times10^{4}$,
so every large-$t$ pooled point comes from \emph{small} $D\ (\le7$k$)$, which floors early and
over-trains --- its steep large-$t$ decay is not representative of a large-$D$ run; (ii) a big
foundation run therefore lives in the (high-$D$, high-$t$) corner that the grid does not cover,
so its centre is an \emph{extrapolation} of the high-$D$ row's decay ($\approx$2, 1.4,
$0.9\times10^{-3}$ at $t=$30, 50, 100k for per-process $D\!\approx\!70$k), and the sweep both
centres the search and \emph{measures} that corner. (The earlier ``is $\vlnr$ still dropping
vs.\ floored?'' rule of thumb is dropped: it was never a measured axis --- only $t$ and $D$
were.)
\end{remark}

% --------------------------------------------------------------------------
\subsection{Which knobs matter, and where the optima land}
% --------------------------------------------------------------------------

Ranking the hyperparameters we still tune by the mean $\abs{\text{Spearman}(\text{HP},\vlnr)}$
within a sweep, $\eta$ is the dominant knob ($0.29$), followed by
\code{regularization\_lambda} ($0.23$), \code{warmup\_frac} ($0.22$), \code{ema\_decay}
($0.20$), and \code{cosanneal\_eta\_min} ($0.20$). (The sampler knobs are excluded --- the
balanced sampler is no longer used, \S\ref{sec:loss}.) For finetune sweeps \code{lr\_scale}
dominates far more strongly ($0.48$), then \code{layer\_decay} and \code{warmup} ($0.29$). The
best finetune learning rate is $\approx$ the pretrain rate: \code{lr\_scale} optimum $\approx
1.0$ (median $1.26$, p5--p95 $0.13$--$7.7$), and \code{layer\_decay} $\approx 0.88$ (p5--p95
$0.74$--$1.0$).

Table~\ref{tab:searchspace} contrasts where the optima actually land against the previously
declared (much wider) ranges; the right column is the redesigned default search space that
these results impose (\S\ref{sec:principles}).

\begin{table}[H]
\centering
\caption{Where HP optima land vs.\ the declared range, and the redesigned default. The
$\eta$ window is centred on Eq.~\eqref{eq:invU} and swept $\pm\tfrac12$ decade
($\eta_{\mathrm{lo/hi}}=\eta_c/3,\ 3\eta_c$ with
$\eta_c(t)=10^{-2}\min[(t/3000)^{0.5},(t/3000)^{-0.55}]$).}
\label{tab:searchspace}
\begin{tabular}{@{}llll@{}}
\toprule
Hyperparameter & Declared (old) & Optima p5--p95 & Redesigned default \\
\midrule
\code{training.lr} & $3.2\text{e-}5$--$3\text{e-}1$ (4 dec) & $3.5\text{e-}4$--$2.1\text{e-}2$ & centre on $\eta_c(t)$, $\pm\tfrac12$ dec \\
\code{regularization\_lambda} & $1\text{e-}11$--$1\text{e-}2$ (9 dec) & $2.5\text{e-}11$--$9.8\text{e-}7$ & $[1\text{e-}10,\,1\text{e-}6]$ \\
\code{cosanneal\_eta\_min} & $1\text{e-}11$--$1\text{e-}6$ & spans full, weakest & fix $\approx 1\text{e-}8$ \\
\code{cosanneal\_warmup\_frac} & $0$--$0.2$ & $0.007$--$0.2$ (med.\ $0.15$) & $[0.05,\,0.2]$ \\
\code{ema\_decay} & $0.9$--$1.0$ & $0.90$--$1.0$ (med.\ $0.96$) & $[0.9,\,0.999]$ \\
\code{fine\_tune.lr\_scale} & $5\text{e-}3$--$50$ (4 dec) & $0.13$--$7.7$ & $[0.1,\,10]$ (2 dec), centre $1$ \\
\code{fine\_tune.layer\_decay} & $0.65$--$1.0$ & $0.74$--$1.0$ & $[0.75,\,1.0]$ \\
\bottomrule
\end{tabular}
\end{table}

\begin{figure}[H]
  \centering
  \includegraphics[width=\textwidth]{hpo_optima_summary.pdf}
  \caption{\textbf{HPO optima across 422 converged sweeps.} $\eta^\ast$ vs width (flat,
  $\mu$P transfers), vs batch (entangled with horizon --- do not read as causal), and vs
  horizon (the inverted-$U$); the $\eta^\ast$ histogram with the recommended band; and the
  finetune \code{lr\_scale} landing near $1$.}
  \label{fig:hpo}
\end{figure}

\begin{figure}[H]
  \centering
  \includegraphics[width=0.86\textwidth]{lr_vs_iterations_by_regime.pdf}
  \caption{\textbf{The inverted-$U$ law.} Geometric-mean best $\eta$ per exact horizon
  $t$, pooled (left) and split by regime (right): $\eta^\ast$ rises as $\sim t^{+0.5}$ up to
  $\tstar\approx3$k steps, then falls as $\sim t^{-0.55}$.}
  \label{fig:iters}
\end{figure}

\begin{figure}[H]
  \centering
  \includegraphics[width=\textwidth]{lr_2d_disentangled.pdf}
  \caption{\textbf{Disentangling horizon from data size} on the clean 2D grid. (A) $\eta^\ast$
  vs horizon at fixed $D$: an inverted-$U$ at every $D$, peak pinned to $\sim3$k steps
  independent of $D$. (B) $\eta^\ast$ vs $D$ at fixed horizon: weak, $\propto D^{+0.15\ldots0.3}$,
  saturating. (C) learning curves flooring later as $D$ grows --- the peak is not the floor.}
  \label{fig:2d}
\end{figure}

% ==========================================================================
\section{Physical feature levers: what the momenta cannot encode}\label{sec:levers}
% ==========================================================================

A recurring question is whether handing the model a physics-motivated input feature
(Def.~\ref{def:coup}) helps. The organizing result is a single principle, verified across
coupling, external-mass, and internal-mass/resonance experiments (all single-seed at matched
HPs; arms are comparable \emph{within} a test, which shares one standardization scheme, but not
\emph{across} tests).

\begin{principle}[Feed only what the external momenta cannot recover]\label{prin:levers}
A feature helps precisely when the quantity it supplies is \emph{not} a function of the
external four-momenta --- an internal or theory input (coupling $\alpha$, a propagator mass) ---
\emph{and} is supplied in a form the network can actually use (a kinematic$\,\times\,$mass
\emph{interaction}, not a near-constant scalar). It is neutral when the information is already
present in the momenta, or already carried by the per-process amplitude offset.
\end{principle}

The experiments (Table~\ref{tab:levers}) each toggle one feature (off $=$ baseline) at matched
HPs on a targeted process:
\begin{itemize}
  \item \textbf{Coupling order} (\code{coupling\_scalars}: the per-event $\log\alpha$ and the
  order labels). When per-dataset standardization removes the constant amplitude offset and the
  kinematics are identical across couplings, the net cannot infer $\alpha$ at all; supplying it
  then cuts $\vlnr$ by $\sim18\times$ ($0.0573\!\to\!0.0031$). With the offset left intact it is
  neutral ($-6.5\%$); where a per-event shape survives (steep running, wide $\alpha$) it wins
  modestly ($-19\%$). This is exactly Principle~\ref{prin:levers}.
  \item \textbf{External mass} (\code{mass\_from\_momenta}, the invariant mass of a leg system)
  is \emph{redundant} --- the mass \emph{is} in the four-momenta; both arms already sit at
  $\vlnr\!\sim\!2\times10^{-5}$.
  \item \textbf{Internal (propagator) mass.} The physically important case: on a scan of the
  internal $Z$ mass in $ee\!\to\!\mu\mu$ across the resonance, the amplitude has a sharp
  Breit--Wigner peak whose location is set by a mass that never appears on an external leg. A raw
  internal log-mass \emph{scalar fed to the model fails} ($+5\%$: a near-constant the net
  ignores). What works is supplying the missing \emph{kinematic$\,\times\,$mass interaction} ---
  the per-event distance to the pole --- directly to the transformer, which collapses the
  mass-scan $U$-shape by three to four orders of magnitude. Routing the same corrected mass
  through the diagram encoder instead is far weaker ($-8\%$, $U$-shape not flattened): the signal
  has to reach the main transformer, not just the diagram embedding.
\end{itemize}

\paragraph{Resonance vs.\ off-shellness (which we adopt, and why).} Two constructions of that
interaction were tried; they \emph{fill the same input slot} and so are mutually exclusive.
\emph{Resonance} $\sqrt s-M$ is the $s$-channel form (distance of the collision energy to the
pole); \emph{off-shellness} $s_{\mathrm{prop}}-M^2$ is the general per-propagator form, valid for
\emph{any} internal line ($t$-channel, multiple propagators, $s$-channel as a special case). On
the single $s$-channel $Z$ test the resonance form is marginally better ($1.8\times10^{-5}$ vs.\
$9.4\times10^{-5}$, both $\sim\!10^3$--$10^4\times$ below baseline). \textbf{We nonetheless adopt
off-shellness}: the small edge of $\sqrt s-M$ is bought by hard-coding the $s$-channel topology,
whereas $s_{\mathrm{prop}}-M^2$ generalizes to the top and Higgs propagators and to non-$s$-channel
diagrams at negligible cost --- exactly what a foundation model spanning many topologies needs.

\begin{table}[H]
\centering
\caption{Feature-lever A/B on $\vlnr$ (best non-regularized val loss, single seed, matched HPs;
``off'' $=$ feature disabled). Each row is a separate targeted test; scales are comparable
\emph{within} a row but not across rows with different standardization (row~1 is a deliberate
pooled-standardization capability probe that makes $\alpha$ unrecoverable by construction).
Resonance and off-shellness populate the same slot (mutually exclusive constructions).}
\label{tab:levers}
\begin{tabular}{@{}llrrl@{}}
\toprule
Lever & Test setting & off & on & Verdict \\
\midrule
Coupling scalars (capability probe) & pooled std, $ee\,u\bar ug\!g\!g$ & $0.0573$ & $0.0031$ & \textbf{win} $18\times$ \\
Coupling scalars (control) & per-dataset std, same & $0.0049$ & $0.0046$ & neutral \\
Coupling scalars (steep-running) & wide $\alpha$, low $\sqrt s$ & $0.0064$ & $0.0052$ & win $-19\%$ \\
Mass from momenta & near-threshold $t\bar t$ & $2.1\text{e-}5$ & $1.7\text{e-}5$ & redundant \\
Internal-mass scalar & $M_Z$ scan, $ee\,\mu\mu$ & $0.0756$ & $0.0795$ & \textbf{lose} \\
Resonance $\sqrt s-M$ ($s$-channel) & $M_Z$ scan & $0.0819$ & $1.8\text{e-}5$ & \textbf{win} $4500\times$ \\
Off-shellness $s-M^2$ (general, \emph{adopted}) & $M_Z$ scan & $0.0763$ & $9.4\text{e-}5$ & \textbf{win} $810\times$ \\
Diagram scanned-mass & $M_Z$ scan, diagram enc. & $0.0777$ & $0.0718$ & weak $-8\%$ \\
\textbf{Off-shellness (flagship, adopted)} & 44 dsets, full data, 30k it & $0.2676$ & $0.0053$ & \textbf{win} $50\times$ \\
\bottomrule
\end{tabular}
\end{table}

\paragraph{The flagship (the production configuration).} The adopted setup --- off-shellness
$s_{\mathrm{prop}}-M^2$ fed directly to the transformer, with coupling scalars, for the internal
$Z$, top, and Higgs propagators (\code{internal\_mass\_pdgs}$=[23,6,25]$) --- is confirmed in the
well-trained regime that matters: $44$ datasets spanning the four resonances
($ee\!\to\!\mu\mu$ $Z$-scan, $ee\!\to\!4\mu$ $Z$-in-$4\ell$, $ee\!\to\!WWb\bar b$ top, and
$ee\!\to\!\mu\mu\tau\tau$ Higgs), on \emph{full} data with fiducial cuts, at $30$k iterations (not
undertraining noise). Off-shellness lowers $\vlnr$ by $\sim\!50\times$
($0.268\!\to\!0.0053$; Fig.~\ref{fig:levers}) --- the headline feature result, and the reason
off-shellness conditioning is on by default in the big-run recipes.

\begin{figure}[H]
  \centering
  \includegraphics[width=0.7\textwidth]{levers_offshell_ab.pdf}
  \caption{\textbf{Flagship off-shellness A/B.} Validation $\vlnr$ vs.\ training step for the
  off-vs-on arms over the $44$-dataset $Z$/top/Higgs mixture (full data, fiducial cuts, 30k
  iterations). Off-shellness conditioning lowers the achievable loss by $\sim\!50\times$.}
  \label{fig:levers}
\end{figure}

% ==========================================================================
\section{Neural scaling and transfer}\label{sec:scaling}
% ==========================================================================

We fit held-out log-amplitude loss against training compute
$C = 2\,F_{\mathrm{step}}\,t$ (MACs, forward$+$backward; $F_{\mathrm{step}}$ known in closed
form from $\mathrm{bs}$, $n$, $d=16\nh$), so cells of different width sit on a fair axis. The
law is the \textbf{floor-aware} power law
\begin{equation}
  L(C) \;=\; A\,C^{-\alpha_C} + L_\infty,
  \label{eq:scaling}
\end{equation}
with an irreducible loss $L_\infty$; the bare power law $L\approx A\,C^{-\alpha_C}$ is only ever a
diagnostic in the regime far above the floor. We report $\alpha_C$ as the local exponent over the
sampled compute range (Fig.~\ref{fig:phase1scaling} shows the phase-1 joint-pretraining learning
curves at each per-process dataset size $D$). A virtual/Born \emph{ratio} target was also tried as a learning
objective; it does not scale meaningfully and is excluded from what follows.

\begin{principle}[Compute exponent falls with final-state multiplicity]\label{prin:mult}
The compute-scaling exponent $\alpha_C$ decreases with the number of final-state particles:
$\alpha_C\approx1.7$--$2.2$ for $2\!\to\!2$, $\approx1.1$--$1.3$ for $2\!\to\!3$, and
$\approx0.96$ for $2\!\to\!4$ (Table~\ref{tab:scaling}), tracking the theoretical lower bound
$\alpha=4/\mathrm{DOF}$ with $\mathrm{DOF}=3n_{\mathrm{fs}}-4$ (Fig.~\ref{fig:alphamult}):
harder, higher-dimensional targets learn more slowly per FLOP.
\end{principle}

Moreover $\alpha_C$ is \textbf{width-independent}: overlaying the $\nh\in\{4,8,16,32\}$ curves at
each data size $D$, they share a slope, so width does not change the compute-scaling exponent ---
consistent with $\mu$P and meaning there is no separate loss-vs-parameters power law; width enters
only through $C$. The compute-optimal (plateau) budget $C^\ast(D)$ where Eq.~\eqref{eq:scaling}
flattens rises only $\sim5\times$ over a $100\times$ range in per-process data, saturating near
$1.1\times10^{17}$ FLOP by $D\gtrsim3\times10^{4}$.

\begin{table}[H]
\centering
\caption{From-scratch (solo) compute-scaling exponents $\alpha_C$, grouped by final-state
multiplicity, and the theoretical bound $\alpha=4/\mathrm{DOF}$ for that multiplicity.}
\label{tab:scaling}
\begin{tabular}{@{}llrr@{}}
\toprule
Process & Type & $\alpha_C$ & $4/\mathrm{DOF}$ \\
\midrule
$ee\to\gamma\gamma$        & QED $2\!\to\!2$ & $2.21$ & \multirow{4}{*}{$2.00$} \\
$ee\to t\bar t$            & EW  $2\!\to\!2$ & $1.95$ & \\
$ee\to WW$                 & EW  $2\!\to\!2$ & $1.81$ & \\
$ee\to u\bar u$            & QCD $2\!\to\!2$ & $1.74$ & \\
\addlinespace
$ee\to u\bar ug$           & QCD $2\!\to\!3$ & $1.34$ & \multirow{3}{*}{$0.80$} \\
$ee\to \gamma\gamma\gamma$ & QED $2\!\to\!3$ & $1.21$ & \\
$ee\to WWZ$                & EW  $2\!\to\!3$ & $1.15$ & \\
\addlinespace
$ee\to u\bar ugg$          & QCD $2\!\to\!4$ & $0.96$ & $0.50$ \\
\bottomrule
\end{tabular}
\end{table}

\begin{figure}[H]
  \centering
  \includegraphics[width=\textwidth]{phase1_scaling.pdf}
  \caption{\textbf{Phase-1 compute scaling of the joint pretraining} ($n_\mathrm{heads}=16$).
  Best held-out loss vs.\ training compute, one panel per per-process dataset size $D$, with the
  bare power law (dashed) and the floor-aware law Eq.~\eqref{eq:scaling} (solid) overlaid; the
  dotted vertical line marks the largest sampled compute $C^\ast$.}
  \label{fig:phase1scaling}
\end{figure}

\begin{figure}[H]
  \centering
  \includegraphics[width=\textwidth,page=2]{phase1_scaling.pdf}\\
  \includegraphics[width=\textwidth,page=3]{phase1_scaling.pdf}\\
  \includegraphics[width=\textwidth,page=4]{phase1_scaling.pdf}
  \caption{\textbf{Phase-1 per-$D$ detail: $D=10^3$, $10^{3.5}$, $10^4$} (top to bottom). All
  trials (light dots) and per-width envelope fits, vs.\ FLOPs (left) and actual GPU-hours (right),
  for $n_\mathrm{heads}\in\{4,8,16,32\}$.}
  \label{fig:phase1detailA}
\end{figure}

\begin{figure}[H]
  \centering
  \includegraphics[width=\textwidth,page=5]{phase1_scaling.pdf}\\
  \includegraphics[width=\textwidth,page=6]{phase1_scaling.pdf}
  \caption{\textbf{Phase-1 per-$D$ detail (cont.): $D=10^{4.5}$, $10^5$} (top to bottom). Same
  layout as Fig.~\ref{fig:phase1detailA}.}
  \label{fig:phase1detailB}
\end{figure}

\begin{figure}[H]
  \centering
  \includegraphics[width=0.66\textwidth]{alpha_vs_multiplicity_overlay.pdf}
  \caption{\textbf{Compute-scaling exponent vs.\ final-state multiplicity.} Fitted $\alpha_C$
  per process family (dots, dotted lines; one colour per family) against the theoretical lower
  bound $\alpha=4/\mathrm{DOF}$ (grey wedge). $2\!\to\!2$ targets sit near the bound; higher
  multiplicity falls below it. Stars are the fine-tuned (full, layer-decay) virtual targets:
  they match the from-scratch exponent --- the transfer advantage is in the prefactor and floor,
  not the slope (\S\ref{sec:transfer}).}
  \label{fig:alphamult}
\end{figure}

% --------------------------------------------------------------------------
\subsection{Foundation pretraining transfers, and how}\label{sec:transfer}
% --------------------------------------------------------------------------

We fine-tune a foundation model onto two held-out NLO-virtual (one-loop) targets --- the virtual
corrections to $ee\!\to\!u\bar u$ and $ee\!\to\!t\bar t$ --- and compare against training those
targets from scratch. Two foundations are used, and their difference matters for reading the
results:
\begin{itemize}
  \item the \textbf{8-process} foundation \emph{includes} the tree-level (LO) $ee\!\to\!u\bar u$
  and $ee\!\to\!t\bar t$ --- the same processes as the fine-tune targets, one order lower --- and
  is trained long ($\approx\!94$k steps at batch $16384$, $\sim\!1.5\times10^{9}$ events);
  \item the \textbf{25-process} foundation \emph{excludes} both targets entirely, and is a short
  run ($15$k steps at batch $1024$, $\sim\!20$ min, $\sim\!1.2$ epochs) --- $\sim\!6\times$ fewer
  steps and $\sim\!100\times$ fewer events than the 8-process one.
\end{itemize}

\paragraph{The transfer advantage is in the prefactor and the floor, not the slope.} The fitted
compute exponent $\alpha_C$ is \emph{roughly the same} for fine-tuning and from-scratch (both
$\approx\!1.7$--$2.2$ on these $2\!\to\!2$ targets; Fig.~\ref{fig:alphamult}). What fine-tuning
buys is a strictly lower curve at every compute budget and, decisively, \textbf{no plateau}:
solo training flattens (even mildly regresses) once it nears its floor within the reachable
budget, whereas the fine-tuned run keeps descending on the same slope to one-to-two orders lower
loss (Fig.~\ref{fig:transfer}). The effect is consistent, not marginal, and holds even for the
short, target-excluding 25-process foundation --- so it is genuine transfer of reusable structure,
not memorization of the target's LO.

\paragraph{Full layer-decay fine-tuning wins; the methods are not equal.} Comparing full
retraining with per-layer learning-rate decay against LoRA, block freezing, EWC, and output-head
reset (Table~\ref{tab:ftmethods}), the exponents end up similar but the methods are far apart in
\emph{where} they reach a given loss. Full layer-decay is best at every intermediate fidelity,
often $1$--$2$ orders below the next method (e.g.\ $ee\,t\bar t$ at $t=400$: $3.2\times10^{-5}$
vs.\ $5.4\times10^{-4}$, $\sim\!17\times$). Read in wall-time this is a $\sim\!10\times$ speedup at
matched loss (full reaches $3\times10^{-5}$ in $0.07$\,h where the others need $\sim\!0.6$\,h;
Fig.~\ref{fig:ftmethods}). Only at the very top fidelity do all methods converge toward the shared
floor. Full retraining with per-layer decay is therefore the default fine-tune method.

\begin{table}[H]
\centering
\caption{Fine-tune method comparison: best $\vlnr$ on $ee\,t\bar t$ NLO-virtual (100k events,
$\nh=8$) at three fidelities. Full layer-decay retraining wins at every intermediate budget; all
methods approach the shared floor only at the top fidelity.}
\label{tab:ftmethods}
\begin{tabular}{@{}lrrrrr@{}}
\toprule
$t$ (steps) & full (decay) & reset-head & freeze & EWC & LoRA \\
\midrule
$400$   & $\mathbf{3.2\times10^{-5}}$ & $5.4\times10^{-4}$ & $5.4\times10^{-4}$ & $5.9\times10^{-4}$ & $1.1\times10^{-3}$ \\
$4000$  & $\mathbf{5.5\times10^{-7}}$ & $1.5\times10^{-6}$ & $1.6\times10^{-6}$ & $5.7\times10^{-6}$ & $4.8\times10^{-6}$ \\
$40000$ & $\mathbf{2.7\times10^{-9}}$ & $2.7\times10^{-9}$ & $4.8\times10^{-9}$ & $1.9\times10^{-8}$ & $8.9\times10^{-9}$ \\
\bottomrule
\end{tabular}
\end{table}

\begin{figure}[H]
  \centering
  \includegraphics[width=0.49\textwidth]{compute_scan_eeuu_wt.pdf}\hfill
  \includegraphics[width=0.49\textwidth]{compute_scan_eettbar_wt.pdf}
  \caption{\textbf{Transfer: the compute scan.} Fair log-space test MSE vs.\ fine-tune wall-time
  on $ee\,u\bar u$ (left) and $ee\,t\bar t$ (right) NLO-virtual, one panel per target dataset
  size $D\in\{10^3,10^4,10^5,10^6\}$: from-scratch vs.\ fine-tuning from the 8- and 25-process
  foundations, and from the leak-free $416$-process big-run pretrains (raw pt25-era encoding vs.\
  the adopted feature set) plus the LO-only $352$-process variant --- the $416$-set curves
  isolate raw multi-process scale from the encoding improvements, and the LO-only curve asks
  whether NLO pretraining data matters for NLO-virtual transfer. The $1$h-pretrain feature-ladder
  markers ($D{=}10^5$ panel) walk raw $\to$ +one-hots $\to$ +mass/coupling scalars $\to$ full
  adopted config at fixed pretrain compute. (Full compute-axis version:
  \code{plots/compute\_scan\_*\_big.pdf}.)}
  \label{fig:transfer}
\end{figure}

\begin{figure}[H]
  \centering
  \includegraphics[width=0.62\textwidth]{finetune_methods_walltime.pdf}
  \caption{\textbf{Fine-tune methods in wall-time.} The five fine-tune methods; full layer-decay
  reaches a given loss soonest.}
  \label{fig:ftmethods_wt}
\end{figure}

\begin{figure}[H]
  \centering
  \includegraphics[width=0.62\textwidth]{finetune_methods_comparison.pdf}
  \caption{\textbf{Fine-tune methods vs.\ compute.} Held-out loss vs.\ training compute for full
  layer-decay, reset-head, freeze, EWC, and LoRA; full layer-decay is lowest at every intermediate
  budget, the methods converging only near the shared floor.}
  \label{fig:ftmethods}
\end{figure}

\begin{principle}[Physics encoding generalizes zero-shot; one-hot PIDs cannot]\label{prin:zeroshot}
The $8$-D physical-property encoding (Def.~\ref{def:enc}) lets a frozen foundation model predict
entirely held-out processes at MSE on standardized log-amplitudes of $0.22$--$0.51$ --- well below
the trivial constant-predictor baseline of $1.0$ --- whereas a one-hot PID model
\emph{cannot even represent} a process containing an unseen particle (e.g.\ $ee\to t\bar t$: PDG
$\pm6$ absent from the training vocabulary). Encoding particles by their quantum numbers, not their
identity, is what makes transfer to new processes possible.
\end{principle}

% ==========================================================================
\section{Architecture and encoding ablations}\label{sec:arch}
% ==========================================================================

\paragraph{Architecture: LLoCa wins on cost \emph{and} accuracy.} At matched parameter count
($\approx1.6$M) on the real 25-process pipeline, the three maintained $\mu$P models differ
enormously in compute (Table~\ref{tab:arch}, left): L-GATr costs $\sim\!20\times$ the FLOPs/step
of LLoCa and the slim variant $\sim\!1.5\times$. A step-matched comparison would be meaningless
here --- the models do not share compute or wall-time per step --- so the fair tests fix the
\emph{budget}. In both an \textbf{equal-compute} and an \textbf{equal-wall-time} $20$-trial DyHPO
sweep (each competitor's step budget set to match LLoCa's $15$k-step training FLOPs, resp.\
wall-clock), \textbf{LLoCa wins} (Table~\ref{tab:arch}, right): $\sim\!3.3$--$3.7\times$ lower
$\vlnr$ than the slim variant and $\sim\!7$--$9\times$ lower than full L-GATr. LLoCa is therefore
the default not by convention but because it is simultaneously the cheapest and the most accurate
at a fixed budget; the slim variant is the affordable fully-tensorial fallback, full L-GATr is
neither.

\begin{table}[H]
\centering
\caption{The three maintained $\mu$P architectures at $\approx1.6$M params, 25-process pipeline.
\emph{Left:} per-step cost ($\mathrm{bs}=1024$, fwd$+$bwd). \emph{Right:} best $\vlnr$ from
budget-matched $20$-trial sweeps --- competitors' step counts set to match LLoCa's $15$k-step
compute (eq.\ compute) and wall-time (eq.\ time). Lower is better; LLoCa wins both.}
\label{tab:arch}
\begin{tabular}{@{}lrr@{\hskip 2.2em}rr@{}}
\toprule
 & \multicolumn{2}{c}{cost} & \multicolumn{2}{c}{best $\vlnr$} \\
\cmidrule(r){2-3}\cmidrule(l){4-5}
Model & Params & GFLOP/step & eq.\ compute & eq.\ time \\
\midrule
LLoCa (\code{LLOCAMuPTransformer}) & $1{,}607{,}060$ & $48.8$ & $\mathbf{5.98\times10^{-3}}$ & $\mathbf{5.98\times10^{-3}}$ \\
L-GATr slim (\code{MuPLGATrSlim})  & $1{,}569{,}765$ & $74.3$ & $2.00\times10^{-2}$ & $2.21\times10^{-2}$ \\
L-GATr (\code{MuPLGATr})           & $1{,}635{,}736$ & $999.2$ & $5.29\times10^{-2}$ & $4.23\times10^{-2}$ \\
\bottomrule
\end{tabular}
\end{table}

\paragraph{Encoding.} The physical-property encoding (Def.~\ref{def:enc}) admits optional
extra channels: a one-hot \emph{spin} embedding (spin is a label, not a magnitude), a binary
\emph{massless} flag, per-column \emph{standardization}, a \emph{color} one-hot, and an MLP
(vs.\ Linear) embedder. Sweeping these under BO ($20$ DyHPO trials/arm) and confirming over $5$
seeds on the 25-process target, the confirmed win is the mass flag and standardization on top of
one-hot spin ($\vlnr = 1.66\pm0.26\times10^{-3}$; Table~\ref{tab:enc}) --- these carry a
physics-consistent gain, largest on colored final states. Adding the color one-hot ($+12\%$) or
the MLP embedder ($+24\%$) did \emph{not} help \emph{on this short, small-data A/B}, but the gaps
are within a few 5-seed bands and the reasoning that motivated them is sound: color is a label
exactly like spin, and an MLP is strictly more expressive than a Linear embed, so neither has any
reason to hurt asymptotically. We therefore keep both \textbf{on by default} (as in the big-run
recipes) and treat them as \emph{open}, to be re-tested on richer data and longer schedules; we
would drop them only if a rigorous test in that regime showed a real regression.
That re-test has since been run at full scale ($448$ datasets, Sec.~\ref{sec:bigrun}): the color
one-hot \emph{confirms its keep} (removing it costs $\sim\!17\%$ on the converged bulk), while the
MLP embedder does not pull its weight and is dropped together with the diagram encoder.

\begin{table}[H]
\centering
\caption{Encoding ablation on the 25-process target ($\vlnr$, 5-seed mean$\pm$std; BO-tuned per
arm), on a short small-data schedule. \code{combo}$=$one-hot spin$+$mass flag$+$standardization.
The color/MLP deltas are within a few seed bands --- treated as open, not settled (kept on by
default).}
\label{tab:enc}
\begin{tabular}{@{}lrl@{}}
\toprule
Encoding arm & $\vlnr$ (5-seed) & rel.\ to combo \\
\midrule
\code{combo}        & $1.66\times10^{-3}\pm2.6\times10^{-4}$ & $1.00$ \\
\code{combo\_color} & $1.85\times10^{-3}\pm3.2\times10^{-4}$ & $1.12$ (short-run) \\
\code{combo\_mlp}   & $2.06\times10^{-3}\pm5.0\times10^{-4}$ & $1.24$ (short-run) \\
\bottomrule
\end{tabular}
\end{table}

% ==========================================================================
\section{The 448-process big run: what fails to learn, and which levers survive at scale}\label{sec:bigrun}
% ==========================================================================

The first full-scale joint pretraining campaign (\code{bigrun\_baseline}: $448$ datasets from the
\code{scan\_bigrun\_100k} recipe, $10^5$ events/process with fiducial cuts, $8601$ steps at
$\mathrm{bs}=16384$, DyHPO over lr/reg/warmup/EMA, $8$ evaluated trials, best geometric-mean
$\vlnr=5.8\times10^{-4}$ at $\eta=4.1\times10^{-3}$ with EMA) doubles as a diagnostic: which
processes a joint model of this size actually learns, and which of the physics levers of
Secs.~\ref{sec:levers}--\ref{sec:arch} still earn their keep at scale.

\paragraph{Failure taxonomy.} The median per-process $\vlnr$ is $2\times10^{-4}$, but $50/448$
datasets sit above $0.05$ --- \emph{the same ones in every top trial}, so the failures are
systematic, not seed noise:

\begin{table}[H]
\centering
\caption{Big-run failure census (best trial; ``failing'' $=$ per-process $\vlnr>0.05$; a loss of
$\sim\!1$ on the standardized log-amplitude means nothing was learned).}
\label{tab:bigrun_census}
\begin{tabular}{@{}lrrl@{}}
\toprule
Dataset category & failing & total & note \\
\midrule
$\alpha_s$-scan variants (\code{\_\_sXX}) & $0$ & $336$ & all learn, down to $10^{-4}$ \\
$M_Z$-scan variants of $2\!\to\!2$ (\code{\_\_mzXX}) & $3$ & $60$ & only the widest-window \code{\_\_mz00}s \\
$m_t/m_H/M_Z^{4\ell}$-scan families & $32$ & $32$ & whole families plateau at $0.1$--$0.3$ \\
plain processes & $15$ & $20$ & worst: $ee\,d\bar d$, $ee\,b\bar b$ at $\vlnr\!\sim\!1.6$ \\
\bottomrule
\end{tabular}
\end{table}

Two mechanisms explain the census; both were isolated against the boring alternatives (no
train/val distribution mismatch --- stats agree to $2\%$; per-dataset preprocessing stats verified
correct; zero NaN/skip-step events in the logs).

\begin{principle}[Resonance needles: a dataset fails when a resonance is a sliver of its window]\label{prin:needle}
A dataset learns its resonance structure only if the resonance occupies a non-negligible fraction
of the sampled kinematic range. The plain $2\!\to\!2$ sets span $\sqrt s\in[25,1000]$\,GeV, so the
$Z$ pole is $\sim\!0.3\%$ of the range ($\sim\!2.6\%$ of events within $10$\,GeV of it) while the
standardized log-amplitude target jumps to $+3.3$--$3.8\sigma$ there (max $7\sigma$): missing the
needle alone contributes $0.3$--$0.4$ to the per-process MSE, which fully accounts for
$ee\,\mu\mu$/$ee\,\tau\tau$. The proof by contrast: the \emph{same} processes' narrow-window
$M_Z$-scan variants (window centered on the shifted pole) learn to $10^{-3}$--$10^{-4}$, and the
only plain processes that learn ($ee\,WW$, $ee\,ZZ$, $ee\,ZZZ$, $ee\,4\gamma$, \dots) are exactly
those whose threshold or physics excludes the $Z$ pole from the window. The all-failing
$m_t/m_H/M_Z^{4\ell}$ families are the same needle in a different variable: their peaks live in
\emph{internal} pair invariant masses ($m_{\mu\mu}$, $m_{Wb}$, $m_{\tau\tau}$), phase-space
corners that flat sampling barely populates, so the whole family plateaus at the
family-average level regardless of the scanned mass value.
\end{principle}

\begin{principle}[Gradient competition under the geometric mean]\label{prin:gradcomp}
The log-space aggregation (Principle~\ref{prin:geomean}) weights each process's gradient by
$1/\mathrm{MSE}_p$: once the bulk of $448$ processes reaches $10^{-4}$, a converged process
outweighs a laggard at $\mathrm{MSE}\sim1$ by $\sim\!10^4$. Consequences observed in every trial,
EMA on or off: laggards never gain traction ($ee\,d\bar d$ flat at $1.6$ from the first
validation), and fragile \emph{learned} processes get overwritten late in training ---
$ee\,\gamma\gamma$ (no resonance at all) trains down to $7\times10^{-4}$ by step $\sim\!7$k and
then regresses to $0.1$--$0.7$. Re-balancing the \emph{sampler} cannot fix this: it changes
process frequency by $O(1)$ factors against a $10^4$ gradient-scale imbalance (consistent with the
balanced-sampler A/B of Sec.~\ref{sec:loss} finding no gain).
\end{principle}

\paragraph{Feature ablation at scale (one-factor-out arms).} Five single-GPU arms re-ran the best
trial's exact HPs with one feature removed (plus one combined arm), the protocol's cheap-shortcut
direction: an ablation that wins at HPs tuned \emph{for the baseline} wins conservatively. Since
the number of stuck processes itself differs between arms and biases the pooled geometric mean,
the decisive column is the geometric mean over the $390$ processes that converged
($\vlnr\le0.05$) in \emph{all six} configurations.

\begin{table}[H]
\centering
\caption{Big-run feature arms at the baseline's best HPs (single seed, $8601$ steps). ``GM
conv.''\ $=$ geometric-mean $\vlnr$ over the $390$-process everywhere-converged subset (the
stuck-dataset--robust comparison); median is over all $448$. Per-process tail values are
run-to-run flaky (e.g.\ $ee\,d\bar d$: $1.63/0.62/1.34$ across baseline/no-diagram/combined) ---
only the aggregates are trusted.}
\label{tab:bigrun_arms}
\begin{tabular}{@{}lrrrrl@{}}
\toprule
Arm & GM all & GM conv. & median & $n_{>0.05}$ & verdict \\
\midrule
baseline (all on)            & $5.79\text{e-}4$ & $2.56\text{e-}4$ & $1.98\text{e-}4$ & $50$ & reference \\
no \code{use\_diagrams}      & $3.88\text{e-}4$ & $1.65\text{e-}4$ & $1.90\text{e-}4$ & $49$ & \textbf{drop diagrams} ($-35\%$, all views) \\
no MLP embedder              & $4.86\text{e-}4$ & $2.21\text{e-}4$ & $3.13\text{e-}4$ & $45$ & tail-driven win; median worse \\
no \code{color\_onehot}      & $6.81\text{e-}4$ & $3.00\text{e-}4$ & $2.47\text{e-}4$ & $57$ & \textbf{keep color} ($+17\%$ if removed) \\
no \code{offshell\_per\_event} & $8.63\text{e-}4$ & $3.74\text{e-}4$ & $4.55\text{e-}4$ & $48$ & \textbf{keep off-shellness} ($+46\%$) \\
no diagrams $+$ no MLP       & $\mathbf{3.85\text{e-}4}$ & $\mathbf{1.59\text{e-}4}$ & $\mathbf{1.88\text{e-}4}$ & $51$ & \textbf{adopted candidate} \\
\bottomrule
\end{tabular}
\end{table}

Readings. (i)~\textbf{Off-shellness reconfirms its flagship status} (Sec.~\ref{sec:levers}) in the
$448$-process regime, and the damage from removing it localizes exactly where the physics says:
the internal-resonance families ($m_t/m_H/M_Z^{4\ell}$ group GM $0.156\!\to\!0.454$;
$ee\,ZH$ $0.08\!\to\!1.32$). (ii)~\textbf{Color one-hot flips from ``open'' to ``keep''}: the
short-run $+12\%$ of Table~\ref{tab:enc} was undertraining noise; at scale removing it hurts
everywhere. (iii)~\textbf{The diagram encoder actively hurts} at scale ($-35\%$ from removing it,
robust in every view and at both failure thresholds), consistent with the earlier finding that
diagram-routed signals are weak (Table~\ref{tab:levers}); (iv)~the \textbf{MLP embedder} is not
load-bearing --- its solo win is tail-driven (median worse) --- and dropping it \emph{with} the
diagrams costs nothing and cures that median regression, so the combined
no-diagrams$+$linear-embed configuration is the adopted candidate: best or tied-best in
every aggregate view while being the simplest and cheapest of the six.

\paragraph{Open threads.} Two structural problems cap what any feature set can do about the
$\sim\!50$ stuck processes; both are actionable and unresolved:
\begin{itemize}
  \item \textbf{Resonance-aware generation} (data-side fix for Principle~\ref{prin:needle}).
  Uniform-$\sqrt s$/flat-phase-space sampling starves the pole regions. Options, in increasing
  invasiveness: drop or window the redundant plain wide-range $2\!\to\!2$ sets (their physics is
  covered by the scan variants); importance-sample $\sqrt s$ around poles at generation time; for
  the internal-invariant-mass needles ($4\ell$, $WWb\bar b$, $\mu\mu\tau\tau$), resonance-aware
  phase-space sampling (\`a la MadGraph channel mapping) is the only real cure.
  \item \textbf{Taming the $1/\mathrm{MSE}$ gradient weight} (optimizer-side fix for
  Principle~\ref{prin:gradcomp}). Candidate one-liner: aggregate
  $\frac1P\sum_p\log(\mathrm{MSE}_p+\tau)$ with $\tau\sim10^{-4}$, which caps a converged
  process's weight at $1/\tau$ while leaving processes above $\tau$ on the exact geometric-mean
  behavior (comparison metrics unchanged). Untested; the decisive check is whether
  $ee\,\gamma\gamma$'s late-training collapse disappears at otherwise-identical HPs.
\end{itemize}

\paragraph{Pairwise attention bias: the diagram signal at the pair level (tested, not adopted).}
Motivated by the delivery-mechanism reading of the diagram results (the off-shellness signal wins
when it reaches the main transformer, Sec.~\ref{sec:levers}) and by the external evidence that
\emph{static} diagram conditioning buys nothing in-distribution even in its home setting
(arXiv:2606.23791: prefix-token diagram embeddings move held-out-process generalization and
$500\times$ fine-tune data efficiency, but only $\sim\!0.01$ AUC on training processes), the
pair-level generalization of \code{offshell\_per\_event} was built and tested
(\code{model.use\_pair\_bias}): every internal propagator's per-event $s-m^2$ (signed-log,
first-batch standardized; propagator identity features log-mass$+$massless flag in v2) drives a
zero-init per-head MLP whose output is \emph{added to the attention logits} of the particle pairs
inside the propagator's resonating cluster --- the pole variable placed where the pair
correlation lives. Implementation: padded-SDPA attention branch (the xformers block-diagonal mask
cannot carry a tensor bias; $\sim\!25$--$30\%$ step-time cost when on), exact-baseline at $t_0$ by
zero-init, CPU equivalence guards in \code{tests/test\_amp.py}.

Three single-seed arms at the adopted candidate config (baseline best HPs unless noted):

\begin{center}\begin{tabular}{@{}lrrrr@{}}
\toprule
Arm & GM all & $\alpha_s$-bulk GM & plain GM & $n_{>0.05}$ \\
\midrule
adopted baseline (no bias)          & $\mathbf{3.85\text{e-}4}$ & $\mathbf{1.24\text{e-}4}$ & $0.209$ & $51$ \\
pair-bias v1 (unbounded, 1-D head)  & $5.26\text{e-}4$ & $1.60\text{e-}4$ & $0.193$ & $52$ \\
pair-bias v2 ($+$identity, cap 8)   & $5.74\text{e-}4$ & $1.89\text{e-}4$ & $\mathbf{0.097}$ & $\mathbf{48}$ \\
pair-bias v2 at $\eta/2$ (lr probe) & $7.37\text{e-}4$ & $2.57\text{e-}4$ & $0.215$ & $53$ \\
\bottomrule
\end{tabular}\end{center}

The signature is consistent and two-sided. (i)~\emph{The mechanism works on its targets}: every
pair-bias arm improves the stuck-tail population (v2: plain GM halved, fewest stuck processes of
any big-run arm), with per-run categorical rescues of exactly the predicted physics ---
$ee\,\text{bhabha}$ $0.81\!\to\!0.012$ and $ee\,\bar\nu\nu$ $0.39\!\to\!0.10$ in v1 (t-channel
poles), $ee\,ZH$ $0.31\!\to\!0.0009$ and $ee\,\gamma\gamma$ $0.41\!\to\!0.009$ in v2 --- but
\emph{which} process gets rescued is run-dependent (the known tail flakiness). (ii)~\emph{The
converged bulk pays a robust $30$--$50\%$ tax} that neither the logit cap, nor propagator-identity
features, nor the dominant HP (the v1 arm's learned bias drifted to $-17$-logit hard masks; the
$\eta/2$ probe made bulk and tail worse simultaneously) removes --- so under the $448$-process
geometric mean the aggregate verdict is a loss, and the feature is \textbf{not adopted}; the
planned fair sweep was cancelled by the lr-probe criterion (bulk did not recover
$\Rightarrow$ the tax is design-intrinsic, not HP-tunable). The lesson joins
Principle~\ref{prin:gradcomp}: a single globally-shared bias function cannot serve $448$
processes whose gradient scales differ by $10^4$ --- the converged bulk's weights dominate the
head's training signal yet any bias motion perturbs their converged attention. The natural
successor couples this to the aggregation thread above: a \emph{per-process gate} on the bias
(active only for processes the loss flags as stuck, or trained under a $\tau$-floored
aggregation that stops the bulk from vetoing it) would let the demonstrated tail rescues stand
without taxing the bulk. Code and arms retained: \code{model.use\_pair\_bias} (default off),
\code{compare\_models/\_bigrun\_arm\_pairbias*}.

% ==========================================================================
\section{Behaviour near amplitude divergences}\label{sec:divergences}
% ==========================================================================

An aggregate $\vlnr$ can be excellent while a thin, high-$|\mathcal{M}|^2$ tail is
mispredicted --- and it is precisely the near-divergent regions (soft/collinear limits,
$s$-channel resonances, thresholds) that dominate cross sections and are hardest for a
surrogate. A useful amplitude model must be accurate \emph{where the physics lives}. We
therefore measure the error in $\log|\mathcal{M}|^2$ \emph{phase-space-locally}, as a
function of the kinematics and of the amplitude magnitude, rather than as a single pooled
number.

\subsection{Two-body phase space and the collinear limit}

For a $2\to2$ process the physical coordinates are the invariant $\sqrt{s}$ and the
c.o.m.\ scattering angle $\cos\theta^*$. Across the jointly-pretrained tree processes of
Table~\ref{tab:div_pretrain} and the fine-tuned NLO-virtual $u\bar u,\,t\bar t$, the model
reproduces $\log|\mathcal{M}|^2$ over the whole plane to MSE
$\Delta\log|\mathcal{M}|^2\sim10^{-7}$, with the residual concentrated at the phase-space
\emph{boundaries} ($\cos\theta^*\to\pm1$, $\sqrt{s}$ extremes) rather than at the physical
peaks.

A genuine divergence occupies a razor-thin sliver: the $t/u$-channel collinear pole of
$e^+e^-\to\gamma\gamma$ ($|\mathcal{M}|^2\sim1/\sin^2\theta$) puts only $0.003\%$ of events
above $\log|\mathcal{M}|^2=6$, all at $|\cos\theta^*|>0.9999$ (peak $9.86$), so uniform
$\cos\theta^*$ binning averages it away. Resolving it needs a coordinate that stretches the
beam directions,
\begin{equation}
  \eta \;=\; \mathrm{sign}(\cos\theta^*)\,\log_{10}\frac{1}{1-|\cos\theta^*|},
  \label{eq:collinear}
\end{equation}
under which the pole becomes a linear ramp, $\log|\mathcal{M}|^2\sim-\ln(1-|\cos\theta^*|)$.
The model follows this ramp up to $\log|\mathcal{M}|^2\approx10$ for $\gamma\gamma$
(Fig.~\ref{fig:div_gg}); $W^+W^-$ gives an asymmetric forward-only ramp ($t$-channel $\nu$
exchange) and $u\bar u$ a flat control, each reproduced with error rising only in the
sparsest, most collinear bins.

\begin{table}[H]
\centering
\caption{Jointly-pretrained (zero-shot) $2\to2$ processes and their angular structure.}
\label{tab:div_pretrain}
\begin{tabular}{@{}lll@{}}
\hline
process & test MSE & angular structure \\
\hline
$e^+e^-\to\gamma\gamma$ & $9.0\mathrm{e}{-}10$ & $t/u$-channel collinear pole at
  $\cos\theta^*\to\pm1$ ($\sim1/\sin^2\theta$), $\sqrt{s}$-independent \\
$e^+e^-\to W^+W^-$ & $2.5\mathrm{e}{-}8$ & $t$-channel $\nu$ exchange, forward peak growing
  with $\sqrt{s}$ \\
$e^+e^-\to u\bar u$ & $6.1\mathrm{e}{-}9$ & $s$-channel $\sim(1+\cos^2\theta)$, no angular
  divergence (control) \\
\hline
\end{tabular}
\end{table}

\begin{figure}[H]
  \centering
  \includegraphics[width=\textwidth]{phase_space_pretrain_aa.pdf}
  \caption{\textbf{$e^+e^-\to\gamma\gamma$, jointly pretrained (zero-shot).}
  $\langle\log|\mathcal{M}|^2\rangle$ over $(\sqrt{s},\cos\theta^*)$ --- truth, model,
  error (top); $\sqrt{s}$/$\cos\theta^*$ projections and predicted-vs-true (middle); and
  the collinear-resolved ramp in $\eta$ of Eq.~\eqref{eq:collinear} (bottom), where the
  model tracks the collinear pole to $\log|\mathcal{M}|^2\approx10$. Error concentrates at
  the boundaries, not the pole.}
  \label{fig:div_gg}
\end{figure}

\subsection{Real infrared singularities}

The decisive test is real gluon emission. For $e^+e^-\to u\bar u g\,(gg)$ the soft
($E_g\to0$) and collinear ($g\parallel$ quark) singularities are exposed by two invariant
resolution variables, the gluon energy fraction and the minimum coloured-pair invariant,
\begin{equation}
  x_g \;=\; \frac{2\,p_g\!\cdot\!Q}{s}, \qquad
  y_{\min} \;=\; \min_{i<j}\frac{(p_i+p_j)^2}{s}
  \qquad (Q=p_{e^-}+p_{e^+}),
  \label{eq:ir}
\end{equation}
the minimum taken over gluon-involving coloured pairs; both $\to0$ in the soft/collinear
limit.

\paragraph{What these variables are, and why they are the right ones.}
Soft and collinear are \emph{distinct} infrared limits. A gluon is \emph{soft} when its
energy vanishes ($E_g\to0$) at any angle --- the eikonal factor gives
$|\mathcal{M}|^2\sim1/E_g^2$; two massless partons are \emph{collinear} when they become
parallel ($\theta_{ij}\to0$) at any energy --- an internal propagator $1/(p_i+p_j)^2$ goes
on shell. Neither limit implies the other, so there are three singular regions: soft-only,
collinear-only, and the soft\,$+$\,collinear overlap where both happen at once and the
divergence is strongest. The gluon energy fraction $x_g=2E_g/\sqrt s$ isolates the
\emph{soft} axis (angle-independent). The pair invariant
$y_{ij}=(p_i+p_j)^2/s=2E_iE_j(1-\cos\theta_{ij})/s$ --- the JADE jet-resolution measure
(JADE, Z.\ Phys.\ C\textbf{33} (1986) 23; Durham/$k_T$ variant, Catani et al., Phys.\
Lett.\ B\textbf{269} (1991) 432) --- vanishes for \emph{either} a soft energy \emph{or} a
collinear angle, so $y_{\min}$ is a single ``distance to the nearest singularity'': one
scalar measuring IR depth, at the cost of conflating the two limits (disentangled by
pairing it with $x_g$, the top rows of Figs.~\ref{fig:div_dalitz}\,ff.).

For a three-body final state these are not arbitrary coordinates. The five-dimensional
phase space splits into \emph{two} shape degrees of freedom --- the energy fractions
$x_i=2E_i/\sqrt s$, constrained by $x_q+x_{\bar q}+x_g=2$, so $(x_q,x_{\bar q})$ is a
complete 2D chart --- and \emph{three} orientation (Euler) angles fixing the event plane
relative to the beam. Every soft and collinear singularity lives entirely in the two shape
variables; the three orientation angles carry only a smooth, pole-free production-angle
modulation that the Dalitz projection integrates over. So the Dalitz plane is the complete
map of the \emph{singular} dynamics --- not literally lossless for orientation-differential
observables, but exhaustive for the divergence structure. On it the leading-order QCD
``antenna'' $|\mathcal{M}|^2\propto(x_q^2+x_{\bar q}^2)/[(1-x_q)(1-x_{\bar q})]$ (Ellis,
Stirling \& Webber, \emph{QCD and Collider Physics}, CUP 1996; Peskin \& Schroeder,
\emph{An Introduction to QFT}, ch.~17) puts the collinear poles on the $x_q\to1$ and
$x_{\bar q}\to1$ edges --- where $1-x_q=(p_{\bar q}+p_g)^2/s$ is the $\bar q g$ invariant
--- meeting at the soft-gluon $(1,1)$ corner. The colour factor is \emph{universal}:
$\gamma$ vs $Z$ production only rescales the normalization and adds smooth orientation
dependence, not the pole structure. This universality is precisely QCD factorization ---
near any soft/collinear limit a real-emission amplitude factorizes into a hard part times a
universal singular factor depending only on the soft energy and the collinear invariant,
i.e.\ exactly $x_g$ and the $y_{ij}$ (Catani \& Seymour, Nucl.\ Phys.\ B\textbf{485} (1997)
291; first NLO 3-jet calculation, Ellis, Ross \& Terrano, Nucl.\ Phys.\ B\textbf{178}
(1981) 421; IR-safe resolution, Sterman \& Weinberg, Phys.\ Rev.\ Lett.\ \textbf{39} (1977)
1436). Because the RAMBO generator samples phase space flat (Kleiss, Stirling \& Ellis,
Comput.\ Phys.\ Commun.\ \textbf{40} (1986) 359), these thin singular slivers are
data-starved by construction --- the premise of the hold-out study of
Sec.~\ref{sec:divergences-holdout}.

\paragraph{Result.} For $u\bar u g$ the model reproduces the soft ($\sim1/x_g^2$) and
soft$+$collinear ($\sim1/y_{\min}$) ramps --- the analytic leading-power slopes ($-2\ln10$
and $-\ln10$ per decade of $\log_{10}$) lie on the measured mean ramps into the deep IR ---
and the full bremsstrahlung Dalitz plane, on which the analytic $C_F$-antenna map is
reproduced to Pearson correlation $0.995$ (Fig.~\ref{fig:div_dalitz}), matching truth
across $\sim\!30$ orders of magnitude in $|\mathcal{M}|^2$ (MSE $5.3\times10^{-5}$). The six-particle $u\bar u gg$, the hardest process in the
mixture, still tracks the mean ramps but is the one case where accuracy \emph{degrades
into} the singular region: the per-event error grows from
$\langle|\Delta\log|\mathcal{M}|^2|\rangle\approx0.03$ at $y_{\min}\sim0.1$ to $0.27$ at
$y_{\min}\sim10^{-4}$ (Fig.~\ref{fig:div_uugg}).

\begin{figure}[H]
  \centering
  \includegraphics[width=\textwidth]{ir_pretrain_uug_dalitz.pdf}
  \caption{\textbf{$e^+e^-\to u\bar u g$ Dalitz plane}, $x_i=2E_i/\sqrt{s}$. Panels:
  analytic LO $C_F$-antenna, truth, model, absolute error. The collinear poles lie on the
  $x_q\to1$ and $x_{\bar q}\to1$ edges and meet at the soft-gluon $(1,1)$ corner; white
  lines on the truth/model panels are the analytic iso-$|\mathcal{M}|^2$ contours. The
  model reproduces truth (analytic-vs-truth Pearson $0.995$) with a residual small even
  along the singular edges.}
  \label{fig:div_dalitz}
\end{figure}

\paragraph{Separating the soft and collinear axes.} Because $y_{\min}$ folds the two limits
together, we isolate the pure \emph{collinear} pole by binning $\log|\mathcal{M}|^2$ against
the $\bar q g$ invariant $1-x_q=(p_{\bar q}+p_g)^2/s$ at \emph{fixed, hard} $x_g$, which
switches the soft pole off (Fig.~\ref{fig:div_uug_coll}). In three disjoint $x_g$ bands the
ramp is a straight line of the universal slope $-\ln10$ ($|\mathcal{M}|^2\sim1/(1-x_q)$) and
is \emph{$x_g$-independent} --- the bands are parallel, offset only by the collinear
splitting kernel, the hallmark of collinear factorization --- and the model tracks each band
to $\langle|\Delta\log|\mathcal{M}|^2|\rangle\lesssim10^{-2}$ into the collinear limit.

\begin{figure}[H]
  \centering
  \includegraphics[width=\textwidth]{ir_pretrain_uug_collinear.pdf}
  \caption{\textbf{Pure collinear ramp, $e^+e^-\to u\bar u g$.} $\log|\mathcal{M}|^2$ vs the
  $\bar q g$ collinear invariant $1-x_q=(p_{\bar q}+p_g)^2/s$, in three bands of fixed hard
  gluon energy $x_g$ (soft pole switched off). Left: truth (faint bands) and model
  (dashed$+$markers) fall on the analytic slope $-\ln10$ (green) and stay parallel across
  $x_g$ --- collinear factorization. Right: the model error into the collinear limit, small
  and $x_g$-independent. Complements the soft ($x_g$) ramp, which isolates the other axis.}
  \label{fig:div_uug_coll}
\end{figure}

\begin{figure}[H]
  \centering
  \includegraphics[width=\textwidth]{ir_pretrain_uugg.pdf}
  \caption{\textbf{$e^+e^-\to u\bar u gg$, the hardest process.} $\log|\mathcal{M}|^2$ vs
  the IR resolution variables of Eq.~\eqref{eq:ir}. The mean ramps (bottom, with the
  analytic leading-power slopes $1/y_{\min}$ and $1/x_g^2$ overlaid) are tracked into the
  deep IR, but the model error (blue, and the error map top-right) grows toward small
  $y_{\min}$ --- the sole case of degradation \emph{into} a singularity.}
  \label{fig:div_uugg}
\end{figure}

\begin{principle} The model holds at the amplitude divergences: across all $2\to2$
(including the $\gamma\gamma$ collinear pole) and $e^+e^-\to u\bar u g$, error rises only at
phase-space boundaries and undersampled bins, not at the physical singularities.
Degradation \emph{into} a divergence appears only for the highest-multiplicity process
($u\bar u gg$) in its deepest soft/collinear limit --- the natural next fine-tuning target.
\end{principle}

\subsection{Extrapolation into a held-out divergence: a hold-out / add-back study on
$e^+e^-\to u\bar u g$}
\label{sec:divergences-holdout}

The degradation \emph{into} the singularity seen for the hardest process
(Sec.~\ref{sec:divergences}) raises a sharper question: can a foundation model that
\emph{never saw} $u\bar u g$ predict its soft/collinear divergence when fine-tuned on
data from which that region is \emph{held out} --- and how much in-region data is needed
to recover it? Because the generator samples phase space flat (RAMBO), the divergent
region is data-starved by construction, so this is a controlled data-density experiment.
All tooling is under \code{analysis/divergences/}.

\paragraph{Setup.} The leave-out base is a 22-process foundation --- the 25-dataset
joint set minus the five final-state-gluon processes that share $u\bar u g$'s IR
($u\bar u g,\,u\bar u gg,\,u\bar u ggg,\,d\bar d g,\,b\bar b g$; recipe
\code{pretrain22\_heldout\_uug.yaml}, $200$k train/process, the adopted architecture of
Sec.~\ref{sec:arch}) --- and reproduces every retained process to test MSE
$\lesssim10^{-3}$. The divergence region is $y_{\min}<c$, with $y_{\min}$ the IR
resolution variable of Eq.~\eqref{eq:ir} and $c=3.9\times10^{-2}$ (the $15$th percentile
of $y_{\min}$ over the $10^6$-event $u\bar u g$ sample). Those events split into $850$k
\emph{FAR} ($y_{\min}>c$, always trained), a fixed $75$k held-out \emph{TEST} set inside
the divergence region (never trained), and a $75$k add-back pool. Five fine-tunes take
the base onto FAR${}\cup{}$(fraction $f$ of the add-back pool) for
$f\in\{0,0.05,0.15,0.5,1\}$ (full retrain with layer decay; the finetune HPs of
\code{finetune\_scaling\_virt\_002} transferred at fixed finetune/pretrain lr ratio,
$4000$ steps each). Each model is then forwarded over the \emph{same} held-out TEST
region and we measure MSE $\Delta\log|\mathcal{M}|^2$ inside it (the evaluation
reconstructs each run's frozen input normalization; a $y_{\min}$ alignment guard
confirms the momenta agree to $<10^{-7}$). $f\!=\!0$ is thus pure extrapolation into an
unseen divergence, $f\!=\!1$ the in-support baseline.

\paragraph{Result --- a small in-region data density recovers most of the divergence.}
The held-out-region error falls steeply and \emph{saturates} with $f$
(Fig.~\ref{fig:div_addback}, Table~\ref{tab:div_addback}): pure extrapolation is poor
(MSE $0.71$, and $11$ in the deepest $y_{\min}<10^{-3}$ bin --- there
$\sqrt{\mathrm{MSE}}\!\approx\!3.3$, i.e.\ the amplitude is mispredicted by
$\sim\!e^{3}$), but adding back just \emph{five percent} of the divergent-region pool
cuts the error $\sim\!24\times$ (to MSE $0.030$), after which the curve flattens toward
the in-support baseline (MSE $2.8\times10^{-3}$ at $f\!=\!1$). Every $f$ degrades
monotonically into deeper IR (the per-$y_{\min}$-bin curves, Fig.~\ref{fig:div_addback}
right). The $f\!=\!0$ failure is diagnostic (Fig.~\ref{fig:div_holdout_f0}): the model's
mean $\log|\mathcal{M}|^2$ \emph{plateaus} instead of following the IR ramp, its
predictions saturate against a ceiling in the predicted-vs-true plane, and the error
piles up in the soft\,+\,collinear corner --- the surrogate cannot invent the
divergence's growth from out-of-region data alone. With the region back in support
($f\!=\!1$, Fig.~\ref{fig:div_holdout_f1}) the ramp is tracked, the ceiling is gone, and
the residual is flat and small across the plane.

\begin{table}[H]
\centering
\caption{Held-out divergence-region ($y_{\min}<c$) error vs add-back fraction $f$, for
$e^+e^-\to u\bar u g$ ($N=75$k held-out events).}
\label{tab:div_addback}
\begin{tabular}{@{}lccccc@{}}
\hline
$f$ (add-back fraction) & $0$ & $0.05$ & $0.15$ & $0.5$ & $1$ \\
\hline
MSE $\Delta\log|\mathcal{M}|^2$ (all held-out) & $0.709$ & $0.030$ & $0.011$ & $0.0063$ & $0.0028$ \\
\quad deepest bin $y_{\min}<10^{-3}$ & $11.1$ & $0.55$ & $0.20$ & $0.10$ & $0.035$ \\
\hline
\end{tabular}
\end{table}

\begin{figure}[H]
  \centering
  \includegraphics[width=\textwidth]{addback_curve.pdf}
  \caption{\textbf{Add-back curve, $e^+e^-\to u\bar u g$.} Held-out
  divergence-region ($y_{\min}<c$) error vs the add-back fraction $f$ (leftmost tick
  $f\!=\!0$ is pure extrapolation). Left: MSE and MAE $\Delta\log|\mathcal{M}|^2$ over
  the whole held-out region; right: MSE split by $y_{\min}$ sub-bin (deeper IR toward the
  top). The steep drop already by $f\!=\!0.05$ and the subsequent flattening show a small
  in-region data density recovers most of the divergence.}
  \label{fig:div_addback}
\end{figure}

\begin{figure}[H]
  \centering
  \includegraphics[width=\textwidth]{heldout_resid_f000.pdf}
  \caption{\textbf{Pure extrapolation ($f\!=\!0$), held-out region}, in the IR plane of
  Eq.~\eqref{eq:ir}. The model mean plateaus off the IR ramp (bottom-left, red vs black),
  predictions saturate against a ceiling (predicted-vs-true, bottom-right), and the error
  concentrates in the soft/collinear corner (top-right) --- the model, never having seen
  the region, cannot reproduce the divergence's growth.}
  \label{fig:div_holdout_f0}
\end{figure}

\begin{figure}[H]
  \centering
  \includegraphics[width=\textwidth]{heldout_resid_f100.pdf}
  \caption{\textbf{In-support baseline ($f\!=\!1$), same held-out region.} With the
  divergent region back in the training support the IR ramp is tracked into the deep IR,
  the predicted-vs-true ceiling is gone, and the error is flat and small across the
  plane --- contrast Fig.~\ref{fig:div_holdout_f0}.}
  \label{fig:div_holdout_f1}
\end{figure}

\begin{principle} A flat-sampled surrogate does not extrapolate \emph{into} an unseen
divergence: held out, $e^+e^-\to u\bar u g$'s soft/collinear region is mispredicted by
$\sim\!e^{3}$ in the deepest bin and the model saturates against a ceiling. But the
divergence is \emph{data-recoverable} --- adding back only $5\%$ of the divergent-region
events cuts the MSE $\sim\!24\times$, and the add-back curve then saturates toward the
in-support baseline.
\end{principle}

\paragraph{Follow-ups.} Separate the soft ($x_g$) and collinear ($1-x_q$) axes in the
\emph{hold-out cut} too (the zero-shot diagnostic of Fig.~\ref{fig:div_uug_coll} already
does this for the pretrained model, isolating the collinear pole at fixed hard $x_g$);
extend to $u\bar u gg$ (its own leave-out base); and a deeper-sampling variant
(regenerate $u\bar u g$ with tighter \code{ptj}/\code{drjj} fiducial cuts to probe below
the current IR cutoff, \code{mg5\_pipeline\_final}).

% ==========================================================================
\section{Training throughput}\label{sec:throughput}
% ==========================================================================

All runs are on V100~32GB (Jean-Zay \code{gpu\_p2}, account \code{itg@v100}); A100 is
unavailable to the account, which matters below. A per-step optimization campaign attacked two
distinct regimes, each fast path proven numerically equal to the original by CPU equivalence
guards (\code{tests/test\_amp.py}, Section~0).

\begin{principle}[The step was data-bound before it was compute-bound]\label{prin:dataloader}
At \code{num\_workers=0}, profiling showed dataloading was $69\%$ ($260$\,ms) of a $380$\,ms
step --- $8192$ serial \code{\_\_getitem\_\_} calls plus a collate concatenation on the main
thread with the GPU idle. Moving to \code{num\_workers=2} with pinned memory and a
\code{non\_blocking} host$\to$device copy overlaps the load under compute and takes the step
from $0.356$ to $0.128$\,s/iter ($2.78\times$), landing on the $\sim0.119$\,s compute floor.
This is now the default (\code{num\_workers: 2}); $nw{=}4$ is no better, $nw{=}6$ oversubscribes
the 8 CPUs and hangs.
\end{principle}

The compute floor was then lowered on two fronts. \emph{Vectorizing the on-GPU hot path} ---
per-event mean pooling, the block-diagonal attention mask (built once per forward, not per
block), the per-process loss (a single segment-mean, not a Python loop over processes), the
$L_2/L_1$ regularizer (one fused \code{\_foreach}), and broadcasting per-particle local frames
across heads instead of materializing a per-head copy --- gives a combined $\sim3.1\times$ at
$nw{=}0$ (dominated by the pooling rewrite, $3.08\times$ alone). A \emph{fused Adam}
(\code{fused\_optimizer}, default on) adds $1.34\times$; \emph{deferring host$\leftrightarrow$device
syncs} (collapsing $\sim4$ syncs/step into one fused \code{stack(\dots).tolist()}, with the NaN
guard moved onto the already-synced grad-norm) adds a further $\sim2\%$ --- small precisely
because, after Principle~\ref{prin:dataloader}, the step is no longer GPU-bound. Each vectorized
path is toggleable (\code{LLOCA\_*} env vars) with the original retained for A/B, and each is
\code{allclose}-verified against it.

\begin{remark}[Carried but inert, and deliberately declined]
\code{allow\_tf32} is carried (default on) but is a \textbf{no-op on V100}; it would give
$\sim2\times$ matmul on A100 at $\sim10^{-3}$ precision, so it must be loss-A/B'd on A100 before
it is trusted --- and A100 is unavailable. Two optimizations were declined on purpose:
converting the $L_2$ penalty to decoupled \code{weight\_decay} (it is folded into $\vlnr$,
checkpoint selection, and the HPO objective, and $\lambda$ is a tuned sweep HP --- converting it
would silently break continuity with all existing sweep results), and \code{torch.compile} (the
compute is dominated by the torch-geometric framesnet and xformers attention, neither of which
compiles --- heavy graph breaks for no gain --- and it would prefix \code{state\_dict} keys,
breaking warm-start/fine-tune reloads). A gradient-checkpointing knob
(\code{checkpoint\_blocks}) trades $\sim$one extra forward for a large activation-memory saving
when memory-bound.
\end{remark}

% ==========================================================================
\section{The compute envelope: VRAM and step time}\label{sec:envelope}
% ==========================================================================

Because Jean-Zay bills a fixed GPU-hour budget, each job should run near the hardware ceiling.
Two quantities bound the operating point on the V100: peak VRAM (limiting batch size and width)
and step time (limiting how many updates fit in an allocation). The standard training batch is
therefore the largest that fits, $\mathrm{bs}=16384$ (at $\nh\!\le\!8$ on the 32\,GB V100);
for the many-process joint runs this is also a statistical requirement, not just throughput ---
at $448$ datasets and a per-dataset geometric-mean loss, $\mathrm{bs}=16384$ gives
$\sim\!36$ events/dataset per batch, where $\mathrm{bs}=1024$ would give only $\sim\!2$ and the
per-dataset MSEs (hence their log-mean) would be noise. With batch size $\mathrm{bs}$,
particles per event $n$, and width $\nh$ ($d=16\nh$), peak VRAM is bilinear,
\begin{equation}
  V_{\text{peak}}(\mathrm{bs},n,\nh) \approx a(\nh)\,\mathrm{bs}\,n + b(\nh)\,\mathrm{bs} + C(\nh)
  \quad[\mathrm{MB}],
  \label{eq:vram}
\end{equation}
with the activation coefficient $a\propto\nh$ and the fixed term $C\propto\nh^2$ (parameters
plus the two AdamW moment buffers, $\approx3\times$ params in float32; the EMA shadow adds one
more params-sized copy). Below $\mathrm{bs}=2048$ the widest tested model ($\nh=64$) fits for
all $n\le12$ on 16\,GB; at a practical $\mathrm{bs}=4096$ one fits $\nh=32$ for $n\le8$. Step
time is set by the MLP ($\mathrm{FLOPs}_{\text{MLP}}\propto n\,d^2$ dominates attention's
$n^2 d$ for $n\lesssim12$): it is linear in $n$ once the GPU is saturated
($\mathrm{bs}\gtrsim2048$) and grows as $t\propto\nh^{1.1}$, so throughput scales roughly as
$1/\nh$. The full six-panel benchmark tables are in \code{notes/scaling\_notes.tex}.

% ==========================================================================
\section{Methodological principles}\label{sec:principles}
% ==========================================================================

The results above impose a concrete protocol on how sweeps, HPO, and comparisons are run.

\begin{principle} Compare on $\vlnr$ --- the best \emph{non-regularized} validation loss on
log-amplitudes (Eq.~\ref{eq:vlnr}) --- never the regularized tracked loss, since $\lambda$ is a
tuned HP. Re-base for preprocessing before comparing two sides (Principle~\ref{prin:perproc}).
\end{principle}
\begin{principle} Never re-sweep $\eta$ across width: $\mu$P transfers it. Tune once at the
smallest width and reuse for all larger widths --- up to a $5\times$ cut on a width-scaling grid.
\end{principle}
\begin{principle} Centre the $\eta$ search by evaluating the measured 2D surface
$\eta^\ast(t,D)$ at your \emph{own} $(t,D)$ --- the inverted-$U$ $\eta_c(t)=10^{-2}\min[(t/3000)^{0.5},(t/3000)^{-0.55}]$
(Eq.~\ref{eq:invU}) lifted onto the high-$D$ row of the grid --- and sweep only $\pm\tfrac12$
decade, not a flat 4-decade prior and not the $t$-only $\eta_c(t)$. Every real run sits past the
peak ($t\gg\tstar=3$k), so the centre \emph{decreases} with horizon; a large-$D$ foundation run
lives in the (high-$D$, high-$t$) corner the grid does not cover, so its centre is an
extrapolation of the high-$D$ decay ($\approx$2, 1.4, $0.9\times10^{-3}$ at $t=$30, 50, 100k for
per-process $D\!\approx\!70$k) and the sweep both centres and measures it. Do not hold $\eta$ at
the $\tstar$ peak, and do not discard $D$.
\end{principle}
\begin{principle} Narrow the secondary knobs to where optima actually land
(Table~\ref{tab:searchspace}): cap \code{regularization\_lambda} at $10^{-6}$, fix
\code{cosanneal\_eta\_min}$\approx10^{-8}$, keep \code{warmup\_frac}$\in[0.05,0.2]$, and drop the
exotic sampler variants ($\abs\rho\approx0.1$, noise). Together with the $\eta$ window this
shrinks the search volume $10$--$100\times$, so halve the trial budget at equal resolution or
explore far more densely at equal budget.
\end{principle}
\begin{principle} Finetune from a pretrain-tuned config: \code{lr\_scale}$\in[0.1,10]$ centred at
$1$ (best finetune $\eta\approx$ pretrain $\eta$), \code{layer\_decay}$\in[0.75,1.0]$; same
reg/warmup/$\eta_{\min}$ rules.
\end{principle}
\begin{principle} Seed new sweep cells from the law, not a flat prior: $\eta_c(t)$ is smooth in
$t$ and width-invariant, so a good rate at one horizon predicts its neighbours. Do \emph{not}
assume batch/data-size invariance --- those axes were never cleanly isolated here.
\end{principle}
\begin{principle} Submit multiple sweeps \emph{interleaved} (round-robin), not all-at-once: DyHPO
trials only inform siblings that \code{observe} before they \code{suggest}, so a simultaneous
launch suggests against a stale surrogate (Def.~\ref{def:dyhpo}). When there are not enough
concurrent jobs for interleaving to serialize anything --- a lone sweep, or a moment of abundant
free GPUs --- fall back to \textbf{sequential batches}: split the trials into (by default) three
dependency-chained waves, so each wave \code{observe}s before the next \code{suggest}s and the
single-fidelity optimizer always has data to fit (Rem.~\ref{rem:dyhpo}).
\end{principle}
\begin{principle} A/B fairly and cheaply: reuse the existing baseline (the best of its own short
HPO sweep --- never a single arbitrary run); first try the feature at the baseline's best HPs and
stop if it already wins; only if it does not, sweep the feature with the same HPO and compare
best-vs-best on $\vlnr$, re-basing for any preprocessing difference first.
\end{principle}

\end{document}
